\documentclass[11pt]{article}
\usepackage[margin=1in]{geometry}
\usepackage{booktabs}
\usepackage{fancyvrb}
\usepackage{enumitem}
\usepackage{needspace}
\setlist{itemsep=1pt,topsep=2pt}

\title{Assignment 2}
\author{Roll No.\ 230217}
\date{}

\begin{document}
\maketitle

\section*{Setup}
Everything is built with \texttt{make} and was tested on GPU3 (NVIDIA A40, CUDA 12.6).

%------------------------------------------------------------
\section*{1\quad Problem 1: 3D convolution}
\begin{Verbatim}[frame=single,fontsize=\footnotesize]
make problem1
./problem1/problem1.out   # N: edit Input_Width (1<<6, 1<<7, 1<<8) in 230217-prob1.cu
\end{Verbatim}
\begin{itemize}
\item \textbf{Basic kernel:} one thread per output, $8\times8\times8$ blocks, input
 and filter read from global memory.
\item \textbf{Optimised kernel:} (a) a $12\times12\times12$ shared-memory tile per
 block (the $8^3$ outputs plus a halo of 2, ghost cells written as 0); (b) the
 filter in \texttt{\_\_constant\_\_} memory; (c) the innermost loop unrolled by hand.
\end{itemize}
\texttt{check\_result} reports ``No differences found'' for both kernels at
$N=64,128,256$.

\begin{center}
\begin{tabular}{rrrrrrr}
\toprule
$N$ & CPU & Basic kernel & Opt kernel & Opt vs basic & Basic e2e & Opt e2e \\
\midrule
64  & 38  & 0.047  & 0.026  & 1.80$\times$  & 0.80  & 0.76 \\
128 & 337 & 0.293 & 0.143 & 2.04$\times$ & 4.19 & 4.00 \\
256 & 2754 & 2.321 & 1.106 & 2.10$\times$ & 26.94 & 24.48 \\
\bottomrule
\end{tabular}

\smallskip
{\small Table 1: Problem 1, ms, best of 6 runs (each averages 10 launches); e2e = with copies.}
\end{center}

The optimised kernel is 1.80--2.10$\times$ faster than the basic one, because the
tile reads each input value from global memory once per block instead of up to
125 times. End to end the two are almost equal, since the copies dominate.

%------------------------------------------------------------
\section*{2\quad Problem 2: grid search}
\begin{Verbatim}[frame=single,fontsize=\footnotesize]
make problem2                 # builds problem2-v1.out, -v2.out, -v3.out from 230217-prob2.cu
cd problem2; ./problem2-v1.out   # also v2, v3; reads disp.txt, grid.txt; writes results-vN.txt
\end{Verbatim}
The grid has $13^{10}\approx1.4\times10^{11}$ points. A point is one 64-bit number,
so the matches can be sorted and written in the loop order of
\texttt{problem2-v0.c}.
\begin{itemize}
\item \textbf{(i) Simple port:} the host loops over variables 1--2 (169 launches); one
 thread per combination of variables 3--6 loops over variables 7--10 and
 recomputes all ten rows for every point.
\item \textbf{(ii) Optimised:} the same exhaustive search with running row sums reused
 across the inner variables, an early exit at the first failing row, and only
 variable 1 on the host (13 launches of $13^5=371{,}293$ threads). Shared-memory
 tiling and prefetching are not used: the 150 input numbers are in constant memory.
\item \textbf{(iii) Thrust:} slices of $2^{28}$ points; \texttt{transform} writes a
 pass/fail flag per point, \texttt{count\_if} sizes the output and
 \texttt{copy\_if} collects the matches.
\end{itemize}
All three match \texttt{problem2-v0.c} byte for byte on $6^{10}$ points (2071
results), $8^{10}$ (9316) and the full grid (11,608).

\begin{center}
\begin{tabular}{lrrrr}
\toprule
Version & Best (s) & Worst (s) & vs (i) & vs CPU \\
\midrule
CPU \texttt{problem2-v0.c} & 676.2 & --- & --- & 1$\times$ \\
(i) simple port    & 36.13 & 80.9 & 1.00$\times$ & 18.7$\times$ \\
(ii) optimised     & 6.28 & 14.7 & 5.76$\times$ & 107.7$\times$ \\
(iii) Thrust       & 34.89 & 91.2 & 1.04$\times$ & 19.4$\times$ \\
\bottomrule
\end{tabular}

\smallskip
{\small Table 2: Problem 2, full grid, \texttt{Total time} printed by each program; 10 interleaved runs per version, the best started with an idle GPU, the worst with other jobs running.}
\end{center}

(ii) is 5.76$\times$ faster than (i) (reused sums, early exit, larger launches).
(i) and (iii) are about equal on an idle GPU; I did not profile why. Under
load all times grew by 2.2--2.6$\times$ and (i) and (iii) changed places.

%------------------------------------------------------------
\Needspace{16\baselineskip}
\section*{3\quad Problem 3: counting Bloom filter}
\begin{Verbatim}[frame=single,fontsize=\footnotesize]
make problem3                 # builds problem3-vanilla.out and problem3-blocked.out
cd problem3; ./problem3-vanilla.out N $PWD/insert_100000.bin \
    $PWD/search_100000.bin $PWD/delete_100000.bin      # same for problem3-blocked.out
\end{Verbatim}
The filter is an array of $N$ 8-bit counters and every key maps to $k=3$ counters:
\begin{Verbatim}[frame=single,fontsize=\footnotesize]
insert(key):   for i in 1..3: counter[hash_i(key)] += 1                  # atomicAdd
search(key):   return all(counter[hash_i(key)] > 0 for i in 1..3)
delete(keys):  kernel 1: found[j] = search(keys[j])         # before any decrement
               kernel 2: if found[j]:                      # atomicCAS, never below 0
                           counter[hash_i(keys[j])] -= 1   for i in 1..3
blocked:       hash_i(key)  ->  hash_4(key, N/128)*128 + hash_i(key, 128)       # i = 1..3
\end{Verbatim}
Four counters share one 32-bit word (CUDA has no 8-bit atomics); counter $p$ is
byte $p\%4$ of word $p/4$. One thread handles one key, and duplicates in a
trace are removed on the host first. \texttt{hash1}--\texttt{hash4} are separate
functions returning a value in $[0,n)$. The blocked variant confines a key to
one 128-counter block (one cache line).
\texttt{check\_result()} replays the operations on a host copy of the counters
and a \texttt{std::set} of members, compares every counter after insert and
after delete, and counts false negatives and false positives (false positives
over all searches). Traces come from \texttt{trace\_generator.cpp}
($n=100{,}000$ in the submission); with 9.9 million keys per operation all
checks also pass.

In every run of Table 3, 0 counters differ from the host model and there are 0
false negatives.

\begin{center}
\begin{tabular}{llrrrr}
\toprule
 & & \multicolumn{2}{c}{vanilla} & \multicolumn{2}{c}{blocked} \\
Trace & $N$ & insert & delete & insert & delete \\
\midrule
$n=10^5$       & $10^6$        & 0.547\% & 0.389\% & 0.662\% & 0.554\% \\
$n=2\cdot10^6$ & $2\cdot10^7$  & 0.583\% & 0.390\% & 0.665\% & 0.500\% \\
$n=10^5$       & $3\cdot10^5$  & 8.42\%  & 7.73\%  & 8.52\%  & 7.91\% \\
\bottomrule
\end{tabular}

\smallskip
{\small Table 3: false-positive rate after the insert and after the delete (a third of the searches are for never-inserted keys, two thirds after the delete).}
\end{center}

The blocked variant has more false positives in every row, because keys in the
same block compete for its 128 counters.

\textbf{Deleting never-inserted keys.} With 10,020 never-inserted keys among the
43,354 deleted keys ($N=3\cdot10^5$) the GPU counters still equal the host model
in all 3 runs. 2469 (vanilla) and 2549 (blocked) of those keys looked present
and were decremented, causing 2025 and 1989 false negatives for real members;
no counting Bloom filter can avoid this.

\Needspace{14\baselineskip}
\textbf{Profiling} (Nsight Compute, first launch of each kernel, $n=2\cdot10^6$,
$N=2\cdot10^7$):

\begin{center}
\small
\begin{tabular}{lrrrr}
\toprule
 & \multicolumn{2}{c}{Vanilla} & \multicolumn{2}{c}{Blocked} \\
Metric & insert & search & insert & search \\
\midrule
Kernel duration & 1.46 ms & 625 $\mu$s & 777 $\mu$s & 371 $\mu$s \\
DRAM traffic & 689 MB & 347 MB & 431 MB & 220 MB \\
Memory bandwidth utilisation (DRAM, \% of peak) & 67.9 & 79.8 & 79.8 & 85.5 \\
Arithmetic intensity (FLOP/byte) & 0 & 0 & 0 & 0 \\
\quad integer instructions per DRAM byte & 0.21 & 0.44 & 0.34 & 0.71 \\
Achieved occupancy (\%) & 69.8 & 81.2 & 72.2 & 81.8 \\
Warp efficiency (active threads / 32) & 32/32 & 32/32 & 32/32 & 32/32 \\
\quad threads not predicated off (of 32) & 29.9 & 30.1 & 30.7 & 30.8 \\
L2 cache hit rate (\%) & 9.7 & 10.4 & 36.5 & 19.9 \\
Long-scoreboard stalls (\% of warp-active time) & 22.6 & 50.9 & 22.2 & 57.5 \\
\bottomrule
\end{tabular}

\smallskip
Table 4: Nsight Compute metrics.
\end{center}

There are no floating-point instructions, so the arithmetic intensity is 0
FLOP/byte; 32 threads are active per warp (about 30 not predicated off). Top
three bottlenecks reported by the profiler:
\begin{enumerate}
\item \textbf{Uncoalesced global accesses:} 84\% (insert) and 83\% (search) of the L2
 sectors are excess; search uses only 5.1 of every 32 bytes fetched.
\item \textbf{Memory-bound, long-scoreboard stalls:} DRAM at 68--85\% of peak; a search
 warp waits 288 (vanilla) / 202 (blocked) cycles for loads.
\item \textbf{Low issue rate and occupancy:} an instruction every 35--146 cycles,
 occupancy 70--84\%.
\end{enumerate}
Blocking makes insert 1.9$\times$ and search 1.7$\times$ faster with 36--37\% less
DRAM traffic (the three counters of a key share a cache line), at the price of
more false positives.

\end{document}
